\documentclass[11pt]{article}
\usepackage{amsmath,amssymb,amsthm,hyperref}
\newtheorem{lemma}{Lemma}
\title{Tensor plane tokens: the exact correction interface}
\date{24 November 2026}
\author{Label audit research checkpoint}
\begin{document}
\maketitle

\paragraph{Status and scope.}
This is a written audit, not a Lean theorem. It records a positive joint
construction for the line-rank part, and the first unresolved equations for
the point-correction part. The pole-direction lemma applies only to an exact
identity on an open token box or a sufficiently rich Cartesian token grid.
It does not rule out dictionaries chosen after seeing a finite table, an
approximate strict-sign decoder, or the original unrestricted model.

\section{Tensor incidence and partial diagonals}

Take $c$ affine planes over a field of order $q$. There are
$M=q^{2c}$ tensor points and $M$ products of nonvertical lines. A row point
$p=(p_1,\ldots,p_c)$ is represented by one-hot coordinates in each of $c$
blocks of size $q^2$. Add a one-hot bit block of size $a$ and a one-hot
column block of size $s$. The target domain has $Mas$ inputs in ambient
dimension $cq^2+a+s$.

For slope tuples $\alpha\in\mathbb F_q^c$, choose additive scalar poles
$P_\alpha=\sum_l p_l(\alpha_l)$ and residues $\eta_b(\alpha)$.
Line sums use the independent memory token at every point and column.
If $g_l$ and $p_l$ have different first coordinates for every $l$, the
off-diagonal point contribution is
\[
 \frac{\eta_b(\alpha(p,g))}{P_{\alpha(p,g)}-t_g(j)}.
\]
The denominator is affine in the one-hot block coordinates. The numerator
is an arbitrary tensor in those coordinates and is generally not affine.

Already for $c=2$, if $g_1=p_1$ and the second coordinates have a defined
nonvertical slope $\alpha_2$, the contribution is instead
\[
 \sum_{\alpha_1\in\mathbb F_q}
 \frac{\eta_b(\alpha_1,\alpha_2)}
 {p_1(\alpha_1)+p_2(\alpha_2)-t_g(j)}.
\]
With distinct active poles and nonzero residues this is a genuine
multiple-pole function of the independently chosen token. It cannot be
replaced by the generic single-pole correction as a rational identity.
For a matching coordinate set $S$, the analogous partial diagonal has
$q^{|S|}$ slope terms. Vertical coincidences between unequal points can
be omitted as in the original compiler, or treated by adding the vertical
line class; they do not remove the equal-point partial diagonals.

An arbitrary real residual on tensor rows and columns can be interpolated
with $(q^2+1)^c$ shared monomial heads. But arbitrary bit-dependent residuals
cost $(a+1)(q^2+1)^c$ in that argument. This is roughly $aM$, so it destroys
the intended $a$-bit gain. A successful tensor construction must exploit
additional structure in both the residue tensor and its partial sums.

\section{Positive joint compilation of the line part}

The line-rank contribution is easier than the point correction. Put
$\ell_{L,l}(x)=\sum_{u\in L_l}x_{l,u}$, which is zero or one on a one-hot
block. For every product line $L$ let $Z_L(z)$ be its arbitrary column-linear
line-sum memory. Two heads with numerator $Z_L(z)$ and opposite signs have
denominators
\[
 D+\varepsilon C-\varepsilon\sum_l\ell_{L,l}(x)
 -\varepsilon\sum_b(\eta_b(L)+C)y_b,
 \quad
 D+\varepsilon C-\varepsilon\sum_l\ell_{L,l}(x)
 -\varepsilon C\sum_b y_b.
\]
Take $C>\max_{L,b}|\eta_b(L)|$. Every displayed nonzero slope has the
reverse orientation, and a sufficiently small positive $\varepsilon$
makes all denominators positive on the full Boolean cube. On the bit axis
the pair cancels exactly; on the column axis its numerator is zero.

On a one-hot bit $b$, subtract all row terms missing at least one of the
$c$ row blocks. The remaining full-block coefficient is
\[
 K_{[c]}(\varepsilon,\ldots,\varepsilon;D-\varepsilon\eta_b(L))
 -K_{[c]}(\varepsilon,\ldots,\varepsilon;D)
 =\frac{c!(c+1)}{D^{c+2}}\varepsilon^{c+1}\eta_b(L)
  +O(\varepsilon^{c+2}).
\]
After one common scaling, it converges to the required arbitrary line
residue times the product-line selector. The lower-row corrections across
\emph{all} lines share the reciprocal basis proved in
\texttt{JointCellFilteredKernels.lean}. Their number is at most
\[
 (a+1)\bigl((q^2+1)^c-q^{2c}\bigr).
\]
Thus the whole line part costs at most
\[
 2M+(a+1)\bigl((q^2+1)^c-M\bigr),
\]
which is $O(M)$ when $ac/q^2$ is bounded. All label-dependent $Z_L$ are
arbitrary throughout this calculation. A common small filler on a fixed
weight slice makes zero slopes strict while preserving target values and
paired axes. Since the eventual target set is finite, the limiting
line-part approximation suffices whenever the complete construction has a
strict margin. This does not solve the point correction or assume its
existence.

\section{A sample-dependent freedom that remains admissible}

The scalar pole pencil need not be fixed across memory points. For every
point $g$, choose an arbitrary pencil $P_g(\alpha)$ and an arbitrary token
$t_g(j)$ for each column, with the uniform bounds
\[
 1-\delta<P_g(\alpha)<1,\qquad
 0<t_g(j)<\delta,\qquad (r+s)\delta<1.
\]
The line aggregate can be
\[
 Z_L(j)=\sum_{g\in L}\frac{1}{P_g(\operatorname{slope}L)-t_g(j)}.
\]
For a query point $p$, the weighted pencil identity still gives its own
full decoder plus one off-diagonal term for every other point with a
defined slope. That term has denominator
$P_g(\operatorname{slope}(p,g))-t_g(j)$, so it is still one admissible
one-hot row Cauchy correction. Consequently the exact axis-zero budget
remains $3q^2+4$.

This statement is kernel checked as
\texttt{jointCell\_exists\_affinePlane\_point\_poles} in
\texttt{JointCellSampleDependentPoles.lean}. The point pencils and tokens
may be chosen after seeing the entire finite label table. Its targeted
build passed with 9123 jobs; its transitive axioms are
\texttt{propext}, \texttt{Classical.choice}, and \texttt{Quot.sound}.
The theorem grants this freedom, but asserts no stronger decoder capacity.

The query residues must remain shared across the memory points in this
factorization. In the tensor construction, additive pencils
$P_g(\alpha)=\sum_l p_{g,l}(\alpha_l)$ remain affine in the one-hot row
blocks for each fixed correction point $g$. The arbitrary residue tensor
and partial diagonals are still unresolved.

\section{Vector tokens and exact pole transport}

Let a point memory now be $t_g(j)\in\mathbb R^r$. A vector Cauchy line
kernel would have denominator
\[
 P_\alpha-u_\alpha\cdot t_g(j).
\]
Full-cube positivity can be enforced, for example, by positive coefficients
$u_{\alpha h}$, tokens in $(0,\delta)^r$, and a margin exceeding
$s\delta\sum_hu_{\alpha h}$ together with every row and bit slope budget.
Positive denominators at the target inputs alone do not establish this
condition.

One original affine correction head for point $g$ may store one fixed
linear functional $v_h\cdot t_g(j)$ in its column coefficients. Its
denominator at a query row is
$D_h(\text{row})-v_h\cdot t_g(j)$, with $v_h$ independent of the row.
The forbidden term is a row-dependent coefficient multiplying a column
token coordinate. An exact $H$-head correction would therefore require,
at each relevant row, an identity of the form
\[
 \frac{A}{P-u\cdot t}
 =R(t)+\sum_{h=1}^H
 \frac{A_h(t)}{D_h-v_h\cdot t},
\]
where $A\ne0$, the numerators $A_h$ and $R$ are affine, and the fixed
normal vectors $v_h$ are shared by all rows for that point module.

\begin{lemma}[Unmatched pole directions cannot disappear]
Suppose $u\ne0$ and the displayed identity holds on a nonempty open set on which all
denominators are nonzero. Then the affine hyperplane
$P-u\cdot t=0$ equals one of the hyperplanes $D_h-v_h\cdot t=0$.
In particular $u$ is proportional to some $v_h$.
\end{lemma}
\begin{proof}
Clear denominators. A polynomial zero on a nonempty open set is zero
identically. If no correction denominator is proportional to $P-u\cdot t$,
choose a point on the latter hyperplane outside the finitely many
intersections with correction hyperplanes. Such a point exists because a
finite union of proper affine subspaces does not cover an affine space
over the infinite field $\mathbb R$. At that point the cleared left side
is $A$ times a product of nonzero correction denominators, while every
term on the right contains the zero factor $P-u\cdot t$. This is a
contradiction. The one-dimensional case is the same argument at the pole.
\end{proof}

After clearing denominators the polynomial in this proof has degree at
most $H+2$, allowing the affine numerators and affine $R$. Consequently
the same conclusion holds if the identity is required on a Cartesian
token grid with more than $H+2$ distinct values in each coordinate and no
zero denominators there. Repeated one-variable interpolation makes the
cleared polynomial identically zero. With constant numerators and $R=0$,
more than $H$ values in each coordinate suffice.

Thus $q^c$ pairwise nonproportional active line directions would require
at least $q^c$ correction directions per point in this exact interface.
If all line directions instead lie among $r$ fixed directions, vector
memory separates into $r$ scalar-token Cauchy families. Their line and
partial-diagonal residue compatibility is still required; the existence
of $r$ independent scalar parameters alone proves no extra saving.

\paragraph{What remains possible.}
In an actual table there are only $s$ columns. Tokens can be chosen after
seeing all independently specified labels. They need not form an open set
or a Cartesian grid, and exact lookup values are unnecessary when only
strict signs matter. The pole argument therefore leaves open a
sample-dependent vector dictionary using a small number of heads per
point. Its target is to realize roughly $a\gg H\log n$ independently
specified bits per point and column, while charging the line part,
partial diagonals, both axes, and all earlier signs.

\section{Additive residues and the admissible exterior gap}

An additive residue tensor
\[
 \eta_b(\alpha)=\eta_{b,0}+\sum_l\eta_{b,l}(\alpha_l)
\]
does give an affine numerator on generic row states. Its scalar decoder
is a sum of $q^c$ reciprocals, but lies in a residue space of dimension
$1+c(q-1)$. A dimension-only Chebyshev claim for this space on the whole
positive exterior half-line is false. The exact rational certificate
\texttt{joint-cell-additive-residue-certificate.json}, verified by
\texttt{JointCellAdditiveResidueCertificate.py}, uses two pole blocks
\[
 \{0,1,139\},\qquad\{0,23,20737\}.
\]
It gives an additive residue tensor whose reciprocal sum takes six
strictly alternating signs at six rational points in an interval of
length $5$ centered near $220.6196884$. The five-dimensional derivative
Wronskian has opposite signs at $200$ and $250$. These are exact finite
rational computations. They are not Lean declarations and do not
establish an admissible decoder or an asymptotic gain.

The whole-cube slope budget prevents this particular example from being
used directly. In one plane, fix a memory point $g$, write the slope
poles as $0=p_0\leq p_1\leq\cdots\leq p_{q-1}=M$, and consider only
generic row states whose first coordinate differs from that of $g$.
Each slope occurs at $q-1$ such states. Give these coordinates denominator
slopes $p_\alpha-M$. The other $q$ states can have slopes tending to zero
from below because the generic numerator is required to vanish there.
On a one-hot generic target the base $d+M$ then gives denominator
$d+p_\alpha$. The infimum full-cube margin needed by this extension is
\[
 d>\bigl(q(q-1)-1\bigr)M-(q-1)\sum_\alpha p_\alpha.
\]
The contributions add over tensor blocks, before charging column tokens,
bit slopes, and strict negative fillers. For $q=3$ and poles $0,A,B$ the
contribution is $3B-2A$. The certificate therefore requires a gap above
$62580$, far larger than its sample gaps. For $q=2$ the contribution is
zero, so this particular margin argument does not exclude small-gap
binary-plane kernels. It still leaves every other compiler cost to be
charged.

\section{The binary-plane product-selector equation}

For $q=2$, let $\Delta_w f(d)=f(d-w)-f(d)$. At a fixed memory point in
one plane, the query point has four possible states: the point itself,
its vertical mate, and two generic states with the two defined slopes.
Give the generic states weights $w_0,w_1$, the vertical state weight zero,
and the matching state weight $w_0+w_1$. The full-block reciprocal filter
then uses $\Delta_{w_\alpha}$ at a generic state and
$\Delta_{w_0+w_1}$ at a matching state. The line sum instead uses
$\Delta_{w_0}+\Delta_{w_1}$ at a matching state. The exact identity is
\[
 \Delta_{w_0+w_1}
 =\Delta_{w_0}+\Delta_{w_1}+\Delta_{w_0}\Delta_{w_1}.
\]
Across $c$ blocks, the difference between the product filter and line sum
therefore has one term for every nonempty subset of matching coordinates.
Only the term matching all coordinates is a diagonal decoder. Direct
termwise cancellation retains $2^c-2$ proper partial-diagonal families,
each with independently chosen point memories. No saving follows from
this identity alone. Solving the full interacting residual jointly is
not ruled out.

A different attempted selector jet starts with a row-dependent generic
denominator $D_g(x)-t_g(j)$ and perturbs it by
$-\varepsilon\sum_l h_{g,l}(x)$. Its top mixed coefficient is a product
of the selectors divided by $(D_g(x)-t_g(j))^{c+1}$. However, every lower
coefficient still has the row-dependent denominator. The shared
lower-degree monomial correction theorem does not cover these terms.
Making the original pole variation infinitesimal as well yields a common
denominator, but collapses the desired many-pole scalar decoder at the
same leading order. This is the first missing lemma for that proposal.

Finally, the earlier proper-block-face repair has size $5^c-4^c$ when
$q=2$, compared with $M=4^c$ tensor points. Its ratio $(5/4)^c-1$ is
unbounded. A growing binary tensor family consequently needs a new face
repair as well as a decoder and partial-diagonal treatment.

\end{document}
